\documentclass[12pt,a4paper]{article}
\usepackage[utf8]{inputenc}
\usepackage[indonesian]{babel}
\usepackage{amsmath}
\usepackage{amsfonts}
\usepackage{amssymb}
\usepackage{geometry}
\usepackage{hyperref}

\geometry{margin=2.5cm}

\title{Metode Regula Falsi dalam Komputasi Numerik}
\author{Ahmad Fakhrul Bawani}
\date{}

\begin{document}

\maketitle

\section{Pendahuluan}
Metode Regula Falsi, atau sering disebut metode posisi palsu (\textit{false position method}), adalah metode pencarian akar yang bersifat tertutup (\textit{bracketing method}). Berbeda dengan metode Biseksi yang membagi rentang menjadi dua bagian sama besar, metode ini menggunakan garis lurus yang menghubungkan dua titik ujung selang untuk memperkirakan letak akar.

\section{Konsep dan Penurunan Rumus}
Metode ini didasarkan pada interpolasi linear antara dua titik $[a, f(a)]$ dan $[b, f(b)]$ sedemikian rupa sehingga $f(a) \cdot f(b) < 0$. Garis lurus yang menghubungkan kedua titik tersebut akan memotong sumbu-$x$ pada titik $c$, yang dianggap sebagai hampiran akar.

Persamaan garis yang melewati kedua titik tersebut adalah:
\begin{equation}
\frac{y - f(b)}{x - b} = \frac{f(b) - f(a)}{b - a}
\end{equation}

Dengan mencari titik potong sumbu-$x$ ($y=0$), kita peroleh rumus untuk $c$:
\begin{equation}
c = b - \frac{f(b)(b - a)}{f(b) - f(a)}
\end{equation}

Atau sering ditulis dalam bentuk:
\begin{equation}
c = \frac{a \cdot f(b) - b \cdot f(a)}{f(b) - f(a)}
\end{equation}

\section{Algoritma Iterasi}
\begin{enumerate}
  \item Pilih interval $[a, b]$ sedemikian rupa sehingga $f(a) \cdot f(b) < 0$.
  \item Hitung nilai $c$ menggunakan rumus Regula Falsi.
	\item Evaluasi nilai $f(c)$:
  \begin{itemize}
    \item Jika $f(a) \cdot f(c) < 0$, maka akar berada di antara $[a, c]$, set $b = c$.
    \item Jika $f(a) \cdot f(c) > 0$, maka akar berada di antara $[c, b]$, set $a = c$.
  \end{itemize}
  \item Ulangi langkah di atas hingga $|f(c)| < \epsilon$ atau selisih $|x_{n+1} - x_n| < \epsilon$.
\end{enumerate}

\section{Kelebihan dan Kekurangan}
\subsection{Kelebihan}
\begin{itemize}
	\item Selalu konvergen karena merupakan metode tertutup (pasti menemukan akar selama tebakan awal benar).
  \item Biasanya lebih cepat mencapai konvergensi dibandingkan metode Biseksi karena memanfaatkan nilai fungsi, bukan hanya tanda (+/-).
\end{itemize}

\subsection{Kekurangan}
\begin{itemize}
  \item Bisa mengalami konvergensi yang sangat lambat jika salah satu titik ujung selang tetap ("stagnan") dan tidak bergerak menuju akar.
  \item Lebih lambat dibandingkan metode Newton-Raphson yang menggunakan informasi turunan.
\end{itemize}

\end{document}